\setcounter{section}{26}

  \zwischenueberschrift{Koneksi Levi--Civita pada himpunan terbuka}

Misalkan

\mavergleichskette
{\vergleichskette
{ U
}
{ \subseteq }{ \R^n
}
{  }{
}
{    }{
}
{   }{
}
}
{}{}{}
terbuka dan padanya diberikan suatu struktur Riemann melalui fungsi-fungsi $g_{ij}$. Kita mempunyai fungsi koordinat $x_i$ dan medan vektor konstan terkait $\partial_i$. Suatu koneksi linear pada bundel tangen
\mathl{U \times \R^n}{} dijelaskan oleh turunan

\mavergleichskettedisp
{\vergleichskette
{ \nabla_{\partial_i } \partial_j
}
{ =} { \sum_{k = 1}^n\Gamma_{ij}^k\,\partial_k
}
{ } {
}
{ } {
}
{ } {
}
}
{}{}{}
Berikut ini akan terlihat bahwa pilihan berikut menghasilkan suatu koneksi yang berkaitan erat dengan struktur Riemann.

\inputdefinition
{}
{

Misalkan

\mavergleichskette
{\vergleichskette
{ U
}
{ \subseteq }{ \R^n
}
{  }{
}
{    }{
}
{   }{
}
}
{}{}{}
terbuka dan padanya diberikan suatu struktur Riemann melalui fungsi-fungsi $\left(  g_{ij}  \right)$ dengan matriks invers
\mathl{\left(  g^{k l}  \right)}{.} Fungsi-fungsi bernilai riil yang didefinisikan pada $U$,

\mavergleichskettedisp
{\vergleichskette
{ \Gamma_{ij}^k
}
{ =} { { \frac{   1  }{  2 } } \sum_{l = 1}^n g^{kl}(\partial_ig_{jl}+\partial_jg_{il}-\partial_lg_{ij})
}
{ } {
}
{ } {
}
{ } {
}
}
{}{}{}
disebut
\definitionswort {simbol Christoffel}{}
untuk koneksi Levi--Civita.

}

\inputdefinition
{}
{

Misalkan

\mavergleichskette
{\vergleichskette
{ U
}
{ \subseteq }{ \R^n
}
{  }{
}
{    }{
}
{   }{
}
}
{}{}{}
terbuka dan padanya diberikan suatu struktur Riemann melalui fungsi-fungsi $\left(  g_{ij}  \right)$ dengan matriks invers
\mathl{\left(  g^{k l}  \right)}{.} Pada $U$,
\definitionsverweis {simbol Christoffel}{}{}

\mavergleichskettedisp
{\vergleichskette
{ \Gamma_{ij}^k
}
{ =} { { \frac{   1  }{  2 } } \sum_{l = 1}^n g^{kl}(\partial_ig_{jl}+\partial_jg_{il}-\partial_lg_{ij})
}
{ } {
}
{ } {
}
{ } {
}
}
{}{}{}
menentukan suatu
\definitionsverweis {koneksi linear}{}{}
pada

\mavergleichskette
{\vergleichskette
{ TU
}
{ = }{ U \times \R^n
}
{  }{
}
{    }{
}
{   }{
}
}
{}{}{}
yang disebut
\definitionswort {koneksi Levi--Civita}{}
pada $U$.

}



\inputfaktbeweis
{R^n/Offene Menge/Levi-Civita-Zusammenhang/Trivial/Fakt}
{Lemma}
{}
{

\faktsituation {Misalkan

\mavergleichskette
{\vergleichskette
{ U
}
{ \subseteq }{ \R^n
}
{  }{
}
{    }{
}
{   }{
}
}
{}{}{}
suatu
\definitionsverweis {himpunan bagian terbuka}{}{,}
yang dilengkapi dengan
\definitionsverweis {struktur Riemann}{}{}
terinduksi dari $\R^n$.}
\faktfolgerung {Maka
\definitionsverweis {koneksi Levi--Civita}{}{}
pada $U$
\definitionsverweis {trivial}{}{.}
Secara khusus, berlaku pernyataan-pernyataan berikut.
\aufzaehlungvier{Semua
\definitionsverweis {simbol Christoffel}{}{}
memenuhi

\mavergleichskettedisp
{\vergleichskette
{ \Gamma_{ij}^k
}
{ =} { 0
}
{ } {
}
{ } {
}
{ } {
}
}
{}{}{}
untuk setiap $i,j,k$.
}{Berlaku

\mavergleichskettedisp
{\vergleichskette
{ \nabla_{\partial_i }  \partial_j
}
{ =} { 0
}
{ } {
}
{ } {
}
{ } {
}
}
{}{}{}
untuk
\definitionsverweis {medan vektor standar}{}{}
$\partial_i$.
}{Untuk fungsi-fungsi terdiferensialkan $f_j$ berlaku

\mavergleichskettedisp
{\vergleichskette
{ \nabla_{\partial_i }  \left(  \sum_{j = 1}^n f_j \partial_j  \right)
}
{ =} { \sum_{j = 1}^n \partial_i(f_j) \partial_j
}
{ } {
}
{ } {
}
{ } {
}
}
{}{}{.}
}{Untuk suatu medan vektor $V$ dan suatu medan vektor terdiferensialkan $W$ pada $U$, berlaku

\mavergleichskettedisp
{\vergleichskette
{ ( \nabla_V W)(P)
}
{ =} { \left(  D W   \right)_{ P } \left(   V(P)  \right)
}
{ } {
}
{ } {
}
{ } {
}
}
{}{}{.}
}}
\faktzusatz {}
\faktzusatz {}

}
{

Fungsi-fungsi fundamental Riemann $g_{ij}$ semuanya konstan, sehingga turunan parsialnya yang muncul dalam definisi simbol Christoffel lenyap. Karena itu, menurut
Catatan 25.8,
koneksi tersebut adalah koneksi trivial. Sifat-sifat lainnya mengikuti dari
Latihan 25.11.

}

\inputbeispiel{}
{

Misalkan

\mavergleichskette
{\vergleichskette
{ I
}
{ \subseteq }{ \R
}
{  }{
}
{    }{
}
{   }{
}
}
{}{}{}
suatu
\definitionsverweis {interval riil}{}{}
dan misalkan
\maabbdisp {g} { I } {\R_+
} {}
suatu fungsi positif terdiferensialkan, yang kita tafsirkan sebagai
\definitionsverweis {metrik Riemann}{}{}
pada $I$; lihat
Contoh 16.10.
Satu-satunya
\definitionsverweis {simbol Christoffel}{}{}
untuk koneksi Levi--Civita ialah

\mavergleichskettedisp
{\vergleichskette
{ \Gamma
}
{ =} { { \frac{   1  }{  2 } }  g^{-1} g'
}
{ } {
}
{ } {
}
{ } {
}
}
{}{}{,}
yakni, selain faktor di depannya, merupakan
\definitionsverweis {turunan logaritmik}{}{}
dari $g$. Turunan vertikalnya adalah

\mavergleichskettedisp
{\vergleichskette
{ \nabla_\partial \partial
}
{ =} { { \frac{  g' }{  2g } } \partial
}
{ } {
}
{ } {
}
{ } {
}
}
{}{}{,}
dengan $\partial$ menyatakan turunan pada $I$. Jika diterapkan pada suatu fungsi yang terdiferensialkan secara kontinu
\maabb {f} { I } {\R
} {}
\zusatzklammer {atau, setara dengan itu, pada medan arah $f \partial$} {} {,}
maka menurut
Teorema 25.5
berlaku

\mavergleichskettedisp
{\vergleichskette
{ \nabla_\partial (f)
}
{ =} { \partial f + f  { \frac{  g' }{  2g } }
}
{ =} { f' + f { \frac{  g' }{  2g } }
}
{ } {
}
{ } {
}
}
{}{}{.}

}



\inputfaktbeweis
{Riemannsche Mannigfaltigkeit/Zweidimensional/Levi-Civita-Zusammenhang/Christoffelsymbole/Beziehung/Fakt}
{Lemma}
{}
{

\faktsituation {Misalkan

\mavergleichskette
{\vergleichskette
{ U
}
{ \subseteq }{ \R^2
}
{  }{
}
{    }{
}
{   }{
}
}
{}{}{}
\definitionsverweis {terbuka}{}{}
dan dilengkapi dengan suatu struktur Riemann yang diberikan oleh
\definitionsverweis {fungsi-fungsi terdiferensialkan}{}{}
\maabbdisp {g_{11},g_{12},g_{22}} { U } {\R
} {}
serta
\definitionsverweis {struktur Riemann}{}{}
tersebut mempunyai matriks invers $\left(  g^{kl}  \right)$.}
\faktfolgerung {Maka
\definitionsverweis {simbol Christoffel}{}{}
untuk koneksi Levi--Civita pada $U$ adalah

\mavergleichskettedisphandlinks
{\vergleichskettedisphandlinks
{ \Gamma_{ij}^k
}
{ =} { { \frac{   1  }{  2 } } \left(  g^{k1}(\partial_ig_{j1}+\partial_jg_{i1}-\partial_1 g_{ij})   +  g^{k2}(\partial_ig_{j2}+\partial_jg_{i2}-\partial_2 g_{ij})  \right)
}
{ } {
}
{ } {
}
{ } {
}
}
{}{}{.}
Jika

\mavergleichskette
{\vergleichskette
{ Y
}
{ \subseteq }{ \R^3
}
{  }{
}
{    }{
}
{   }{
}
}
{}{}{}
suatu permukaan berorientasi yang dua kali terdiferensialkan secara kontinu dan dilengkapi dengan struktur yang diinduksi oleh ruang sekitar $\R^3$ sebagai
\definitionsverweis {struktur Riemann}{}{,}
serta
\maabbdisp {\varphi} { V } { Y
} {,}

\mavergleichskette
{\vergleichskette
{ V
}
{ \subseteq }{ \R^2
}
{  }{
}
{    }{
}
{   }{
}
}
{}{}{}
terbuka, merupakan suatu parametrisasi lokal $Y$ yang dua kali terdiferensialkan, maka simbol-simbol Christoffel ini berimpit dengan
\definitionsverweis {simbol Christoffel}{}{}
yang didefinisikan dengan bantuan ruang sekitar.}
\faktzusatz {}
\faktzusatz {}

}
{

Pernyataan pertama langsung mengikuti dari
Definisi 26.1.
Pernyataan kedua mengikuti dari pernyataan pertama dan
Lemma 19.5.

}



\inputfaktbeweis
{Offene Menge/Riemannsche Struktur/Christoffelsymbole/Eigenschaften/Fakt}
{Lemma}
{}
{

\faktsituation {Misalkan

\mavergleichskette
{\vergleichskette
{ U
}
{ \subseteq }{ \R^n
}
{  }{
}
{    }{
}
{   }{
}
}
{}{}{}
terbuka dan padanya diberikan suatu
\definitionsverweis {struktur Riemann}{}{}
melalui fungsi-fungsi $\left(  g_{ij}  \right)$. Maka
\definitionsverweis {simbol Christoffel}{}{}
untuk koneksi Levi--Civita pada $U$ memenuhi

\mavergleichskettedisp
{\vergleichskette
{ \Gamma^k_{ij}
}
{ =} { \Gamma^k_{ji}
}
{ } {
}
{ } {
}
{ } {
}
}
{}{}{.}}
\faktfolgerung {}
\faktzusatz {}
\faktzusatz {}

}
{

Hal ini langsung mengikuti dari definisi dan kesimetrian
\mathl{\left(  g_{ij}  \right)}{.}

}

\inputbeispiel{}
{

Kita melanjutkan
Contoh 18.8,
yakni kita meninjau setengah bidang Poincaré $\mathbb H$ dengan fungsi-fungsi fundamental Riemann

\mavergleichskette
{\vergleichskette
{ g_{11}
}
{ = }{ g_{22}
}
{ = }{ { \frac{   1  }{  y^2 } }
}
{    }{
}
{   }{
}
}
{}{}{}
dan

\mavergleichskette
{\vergleichskette
{ g_{12}
}
{ = }{ 0
}
{  }{
}
{    }{
}
{   }{
}
}
{}{}{.}
Menurut
Lemma 26.5,
dari

\mathdisp {{ \frac{   1  }{  2 } } \left(  g^{k1}(\partial_ig_{j1}+\partial_jg_{i1}-\partial_1 g_{ij})   +  g^{k2}(\partial_ig_{j2}+\partial_jg_{i2}-\partial_2 g_{ij})  \right)} {  }

diperoleh

\mavergleichskettedisp
{\vergleichskette
{ \Gamma^1_{11}
}
{ =} { { \frac{   1  }{  2 } }  \left(  y^2  \partial_1  \left( \frac{   1  }{  y^2 } \right)     \right)
}
{ =} { 0
}
{ } {
}
{ } {}
}
{}{}{,}

\mavergleichskettedisp
{\vergleichskette
{ \Gamma^2_{11}
}
{ =} { { \frac{   1  }{  2 } }  \left(  y^2 \left(  - \partial_2  \left( \frac{   1  }{  y^2 } \right)  \right)      \right)
}
{ =} { { \frac{   1  }{  2 } }  \left(  y^2 \left(  -  \left( \frac{   -2 }{  y^3 } \right)  \right)      \right)
}
{ =} { { \frac{   1  }{  y } }
}
{ } {
}
}
{}{}{}
dan dengan cara serupa

\mavergleichskette
{\vergleichskette
{ \Gamma^1_{12}
}
{ = }{ - { \frac{   1  }{  y } }
}
{  }{
}
{    }{
}
{   }{
}
}
{}{}{,}

\mavergleichskette
{\vergleichskette
{ \Gamma^2_{12}
}
{ = }{ 0
}
{  }{
}
{    }{
}
{   }{
}
}
{}{}{,}

\mavergleichskette
{\vergleichskette
{ \Gamma^1_{22}
}
{ = }{ 0
}
{  }{
}
{    }{
}
{   }{
}
}
{}{}{,}

\mavergleichskette
{\vergleichskette
{ \Gamma^2_{22}
}
{ = }{ - { \frac{   1  }{  y } }
}
{  }{
}
{    }{
}
{   }{
}
}
{}{}{.}

}

  \zwischenueberschrift{Koneksi Levi--Civita}

Kita hendak memperluas konsep koneksi Levi--Civita dari suatu himpunan terbuka

\mavergleichskette
{\vergleichskette
{ U
}
{ \subseteq }{ \R^n
}
{  }{
}
{    }{
}
{   }{
}
}
{}{}{}
dengan struktur Riemann ke suatu manifold Riemann sembarang. Kesulitan utamanya adalah menunjukkan bahwa konstruksi langsung pada daerah-daerah peta terbuka saling konsisten pada irisan yang tumpang tindih. Untuk itu, kita memperkenalkan konsep-konsep berikut.

\inputdefinition
{}
{

Misalkan $M$ suatu
\definitionsverweis {manifold terdiferensialkan}{}{}
dan diberikan suatu
\definitionsverweis {koneksi linear}{}{}
pada
\definitionsverweis {bundel tangen}{}{}
$TM$. Koneksi tersebut disebut
\definitionswort {bebas torsi}{,}
jika

\mavergleichskettedisp
{\vergleichskette
{ \nabla_V W- \nabla_W V
}
{ =} { [V,W]
}
{ } {
}
{ } {
}
{ } {
}
}
{}{}{}
untuk setiap
\definitionsverweis {medan vektor}{}{}
$V,W$ yang
\definitionsverweis {terdiferensialkan secara kontinu}{}{}
pada setiap himpunan terbuka.

}

\inputdefinition
{}
{

Misalkan $M$ suatu
\definitionsverweis {manifold Riemann}{}{}
dan diberikan suatu
\definitionsverweis {koneksi linear}{}{}
pada
\definitionsverweis {bundel tangen}{}{}
$TM$. Koneksi tersebut disebut
\definitionswort {metrik}{,}
jika

\mavergleichskettedisp
{\vergleichskette
{ D_V \left\langle  W  ,  Z  \right\rangle
}
{ =} { \left\langle \nabla_V W , Z  \right\rangle  + \left\langle  W  ,  \nabla_VZ  \right\rangle
}
{ } {
}
{ } {
}
{ } {
}
}
{}{}{.}
untuk setiap
\definitionsverweis {medan vektor}{}{}
$V,W,Z$ yang
\definitionsverweis {terdiferensialkan secara kontinu}{}{}
pada setiap himpunan terbuka.

}



\inputfaktbeweis
{Riemannsche Mannigfaltigkeit/Tangentialbündel/Linearer Zusammenhang/Metrisch und torsionsfrei/Koszul-Eigenschaft/Eindeutigkeit/Fakt}
{Lemma}
{}
{

\faktsituation {Misalkan $M$ suatu
\definitionsverweis {manifold Riemann}{}{}
dan diberikan suatu
\definitionsverweis {koneksi linear}{}{}
$\nabla$ pada
\definitionsverweis {bundel tangen}{}{}
$TM$ yang
\definitionsverweis {metrik}{}{}
dan
\definitionsverweis {bebas torsi}{}{.}}
\faktfolgerung {Maka untuk setiap
\definitionsverweis {medan vektor}{}{}
$V,W,Z$ yang
\definitionsverweis {terdiferensialkan secara kontinu}{}{}
pada setiap himpunan terbuka berlaku apa yang disebut \stichwort {rumus Koszul} {}

\mavergleichskettedisphandlinks
{\vergleichskettedisphandlinks
{ \left\langle \nabla_V W , Z  \right\rangle
}
{ =} { { \frac{   1  }{  2 } }  \left(  D_V \left\langle  W  ,  Z  \right\rangle +  D_W \left\langle  Z  ,  V  \right\rangle -  D_Z \left\langle  V  ,  W  \right\rangle  - \left\langle  W  , [V,Z]   \right\rangle - \left\langle  Z  , [W,V]   \right\rangle + \left\langle  V  , [Z,W]   \right\rangle   \right)
}
{ } {
}
{ } {
}
{ } {
}
}
{}{}{.}}
\faktzusatz {Secara khusus, paling banyak ada satu koneksi dengan sifat-sifat ini.}
\faktzusatz {}

}
{

Karena koneksi bebas torsi, berlaku

\mavergleichskettedisp
{\vergleichskette
{ \nabla_VW- \nabla_WV
}
{ =} { [V,W]
}
{ } {
}
{ } {
}
{ } {
}
}
{}{}{,}

\mavergleichskettedisp
{\vergleichskette
{ \nabla_VZ- \nabla_ZV
}
{ =} { [V,Z]
}
{ } {
}
{ } {
}
{ } {
}
}
{}{}{}
dan

\mavergleichskettedisp
{\vergleichskette
{ \nabla_WZ- \nabla_ZW
}
{ =} { [W,Z]
}
{ } {
}
{ } {
}
{ } {
}
}
{}{}{,}
dan identitas pertama juga dapat ditulis sebagai

\mavergleichskettedisp
{\vergleichskette
{ \nabla_VW + \nabla_WV
}
{ =} { 2 \nabla_VW -[V,W]
}
{ } {
}
{ } {
}
{ } {
}
}
{}{}{.}
Karena koneksi metrik, berlaku identitas-identitas

\mavergleichskettedisp
{\vergleichskette
{ D_V \left\langle  W  ,  Z  \right\rangle
}
{ =} { \left\langle \nabla_V W , Z  \right\rangle  + \left\langle  W  ,  \nabla_V Z  \right\rangle
}
{ } {
}
{ } {
}
{ } {
}
}
{}{}{,}

\mavergleichskettedisp
{\vergleichskette
{ D_W \left\langle  Z  ,  V  \right\rangle
}
{ =} { \left\langle \nabla_WZ , V  \right\rangle  + \left\langle  Z  ,  \nabla_WV  \right\rangle
}
{ } {
}
{ } {
}
{ } {
}
}
{}{}{}
dan

\mavergleichskettedisp
{\vergleichskette
{ D_Z \left\langle  V  ,  W  \right\rangle
}
{ =} { \left\langle \nabla_Z V , W  \right\rangle  + \left\langle  V  ,  \nabla_Z W   \right\rangle
}
{ } {
}
{ } {
}
{ } {
}
}
{}{}{.}
Dengan menjumlahkan persamaan-persamaan ini, kita memperoleh

\mavergleichskettedisphandlinks
{\vergleichskettedisphandlinks
{ D_V \left\langle  W  ,  Z  \right\rangle+ D_W \left\langle  Z  ,  V  \right\rangle  - D_Z \left\langle  V  ,  W  \right\rangle
}
{ =} { \left\langle \nabla_V W+\nabla_W V  ,  Z  \right\rangle + \left\langle \nabla_W Z-\nabla_Z W   ,  V  \right\rangle+\left\langle  \nabla_V Z-\nabla_Z V   ,  W  \right\rangle
}
{ } {
}
{ } {
}
{ } {
}
}
{}{}{.}
Kita substitusikan ungkapan-ungkapan di atas ke ruas kanan dan memperoleh

\mavergleichskettealignhandlinks
{\vergleichskettealignhandlinks
{ D_V \left\langle  W  ,  Z  \right\rangle+ D_W \left\langle  Z  ,  V  \right\rangle  - D_Z \left\langle  V  ,  W  \right\rangle
}
{ =} { \left\langle  2 \nabla_VW -[V,W]  ,  Z  \right\rangle + \left\langle  [W,Z]  ,  V  \right\rangle + \left\langle  [V,Z]   ,  W  \right\rangle
}
{ =} { 2\left\langle  \nabla_VW   ,  Z  \right\rangle - \left\langle  [V,W]  ,  Z  \right\rangle + \left\langle  [W,Z]  ,  V  \right\rangle + \left\langle  [V,Z]   ,  W  \right\rangle
}
{ } {
}
{ } {
}
}
{}
{}{.}
Dengan menata ulang persamaan tersebut, diperoleh rumus yang dinyatakan.

Berdasarkan persamaan itu,
\mathl{\left\langle  \nabla_V W , Z  \right\rangle}{} untuk medan-medan vektor sembarang telah ditentukan oleh ruas kanan, yang sama sekali tidak memuat koneksi. Karena hal ini berlaku untuk setiap medan vektor $Z$, turunan vertikal $\nabla_V W$, dan dengan demikian koneksi linearnya, juga ditentukan secara unik.

}

\inputfaktbeweis
{R^n/Offene Menge/Riemannsche Struktur/Zusammenhang aus Christoffel-Symbole/Eigenschaften/Fakt}
{Lemma}
{}
{

\faktsituation {Misalkan

\mavergleichskette
{\vergleichskette
{ U
}
{ \subseteq }{ \R^n
}
{  }{
}
{    }{
}
{   }{
}
}
{}{}{}
\definitionsverweis {terbuka}{}{}
dan padanya diberikan suatu
\definitionsverweis {struktur Riemann}{}{}
melalui fungsi-fungsi $\left(  g_{ij}  \right)$ dengan
\definitionsverweis {matriks invers}{}{}

\mathl{\left(  g^{k l}  \right)}{.} Misalkan $\nabla$ adalah koneksi yang diberikan oleh
\definitionsverweis {simbol-simbol Christoffel}{}{}

\mavergleichskettedisp
{\vergleichskette
{ \Gamma_{ij}^k
}
{ =} { { \frac{   1  }{  2 } } \sum_{l = 1}^n g^{kl}(\partial_ig_{jl}+\partial_jg_{il}-\partial_lg_{ij})
}
{ } {
}
{ } {
}
{ } {
}
}
{}{}{}
sebagai
\definitionsverweis {koneksi Levi--Civita}{}{,}
yang dicirikan oleh

\mavergleichskettedisp
{\vergleichskette
{ \nabla_{\partial_i } \partial_j
}
{ =} { \sum_{k = 1}^n\Gamma_{ij}^k\,\partial_k
}
{ } {
}
{ } {
}
{ } {
}
}
{}{}{.}}
\faktuebergang {Maka $\nabla$ memenuhi sifat-sifat berikut.}
\faktfolgerung {\aufzaehlungvier{Berlaku

\mavergleichskettedisphandlinks
{\vergleichskettedisphandlinks
{ \left\langle  \nabla_{\partial_i }  \partial_j  ,  \partial_k   \right\rangle
}
{ =} { { \frac{   1  }{  2 } }  \left(  D_{\partial_i }  \left\langle  \partial_j  ,  \partial_k   \right\rangle   +D_{\partial_j }  \left\langle  \partial_i  ,  \partial_k   \right\rangle -D_{\partial_k }  \left\langle  \partial_i  ,  \partial_j   \right\rangle   \right)
}
{ } {}
{ } {}
{ } {}
}
{}{}{.}
}{Koneksi ini memenuhi rumus Koszul dari
Lemma 26.10.
}{Koneksi ini
\definitionsverweis {bebas torsi}{}{.}
}{Koneksi ini
\definitionsverweis {metrik}{}{.}
}}
\faktzusatz {}
\faktzusatz {}

}
{

\aufzaehlungvier{Berlaku

\mavergleichskettedisp
{\vergleichskette
{ D_{\partial_k }  \left\langle  \partial_i  ,  \partial_j   \right\rangle
}
{ =} { \partial_k  \left(  g_{ij}  \right)
}
{ } {
}
{ } {
}
{ } {
}
}
{}{}{.}
Selanjutnya,

\mavergleichskettealignhandlinks
{\vergleichskettealignhandlinks
{ \left\langle  \nabla_{\partial_i }  \partial_j  ,  \partial_k   \right\rangle
}
{ =} { \left\langle  \sum_{\ell = 1}^n\Gamma_{ij}^\ell \,\partial_\ell  ,  \partial_k   \right\rangle
}
{ =} { \sum_{\ell = 1}^n\Gamma_{ij}^\ell \left\langle  \,\partial_\ell  ,  \partial_k   \right\rangle
}
{ =} { \sum_{\ell = 1}^n\Gamma_{ij}^\ell g_{\ell k}
}
{ =} { { \frac{   1  }{  2 } } \sum_{\ell = 1}^n  \left(  \sum_{r = 1}^n g^{\ell r}  \left(  \partial_ig_{j r}+\partial_jg_{ir}-\partial_r g_{ij}  \right)   \right)  g_{\ell k}
}
}
{
\vergleichskettefortsetzungalign
{ =} { { \frac{   1  }{  2 } } \sum_{r = 1}^n  \left(  \partial_ig_{j r}+\partial_jg_{ir}-\partial_r g_{ij}  \right)  \left(  \sum_{\ell = 1}^n  g^{\ell r} g_{\ell k}   \right)
}
{ =} { { \frac{   1  }{  2 } }  \left(  \partial_ig_{j k}+\partial_jg_{i k}-\partial_k g_{ij}  \right)
}
{ =} { { \frac{   1  }{  2 } }  \left(  D_{\partial_i }  \left\langle  \partial_j  ,  \partial_k   \right\rangle   +D_{\partial_j }  \left\langle  \partial_i  ,  \partial_k   \right\rangle   -D_{\partial_k }  \left\langle  \partial_i  ,  \partial_j   \right\rangle   \right)
}
{ } {}
}
{}{.}
}{Karena kurung Lie
\mathl{[ \partial_i, \partial_j]}{} trivial pada $U$, menurut bagian (1) rumus Koszul berlaku untuk medan-medan basis $\partial_i$. Untuk kasus umum, lihat
Latihan 26.7.
}{Menurut
Lemma 26.6
dan karena

\mavergleichskette
{\vergleichskette
{ [\partial_i, \partial_j]
}
{ = }{ 0
}
{  }{
}
{    }{
}
{   }{
}
}
{}{}{}
berlaku

\mavergleichskettedisp
{\vergleichskette
{ \nabla_{\partial_i }  \partial_j
}
{ =} { \nabla_{\partial_j }  \partial_i
}
{ } {
}
{ } {
}
{ } {
}
}
{}{}{.}
Pernyataan itu sekarang mengikuti dari
Latihan 26.5.
}{Menurut bagian (1),

\mavergleichskettealignhandlinks
{\vergleichskettealignhandlinks
{ \left\langle  \nabla_{\partial_i }  \partial_j  ,  \partial_k   \right\rangle + \left\langle  \nabla_{\partial_i }  \partial_k  ,  \partial_j   \right\rangle
}
{ =} { { \frac{   1  }{  2 } }  \left(  D_{\partial_i }  \left\langle  \partial_j  ,  \partial_k   \right\rangle   +D_{\partial_j }  \left\langle  \partial_i  ,  \partial_k   \right\rangle   -D_{\partial_k }  \left\langle  \partial_i  ,  \partial_j   \right\rangle    +D_{\partial_i }  \left\langle  \partial_k  ,  \partial_j   \right\rangle +D_{\partial_k }  \left\langle  \partial_i  ,  \partial_j   \right\rangle  - D_{\partial_j }  \left\langle  \partial_i  ,  \partial_k   \right\rangle   \right)
}
{ =} { D_{\partial_i }  \left\langle  \partial_j  ,  \partial_k   \right\rangle
}
{ } {}
{ } {}
}
{}
{}{.}
Pernyataan itu sekarang mengikuti dari
Latihan 26.8.
}

}



\inputfaktbeweis
{Riemannsche Mannigfaltigkeit/Levi-Civita-Zusammenhang/Eindeutige Existenz/Fakt}
{Teorema}
{}
{

\faktsituation {Misalkan $X$ suatu
\definitionsverweis {manifold Riemann}{}{.}
Maka pada
\definitionsverweis {bundel tangen}{}{}
$TX$ terdapat suatu
\definitionsverweis {koneksi linear}{}{}
yang
\definitionsverweis {bebas torsi}{}{,}
\definitionsverweis {metrik}{}{,}
dan ditentukan secara unik.}
\faktfolgerung {}
\faktzusatz {}
\faktzusatz {}

}
{

Menurut
Lemma 26.10,
paling banyak ada satu koneksi semacam itu. Menurut
Lemma 26.11,
pada setiap daerah peta kita dapat memberikan secara eksplisit suatu koneksi yang mempunyai sifat-sifat yang disyaratkan. Karena keunikan yang baru saja disebutkan, koneksi-koneksi yang dikonstruksi dengan cara ini berimpit pada irisan daerah-daerah peta.

}

\inputdefinition
{}
{

Misalkan $X$ suatu
\definitionsverweis {manifold Riemann}{}{.}
\definitionsverweis {Koneksi}{}{}
pada
\definitionsverweis {bundel tangen}{}{}
$TX$ yang didefinisikan secara lokal oleh
\definitionsverweis {simbol-simbol Christoffel}{}{}
disebut
\definitionswort {koneksi Levi--Civita}{.}

}



\inputfaktbeweis
{Riemannsche Mannigfaltigkeit/Levi-Civita-Zusammenhang/Eigenschaften/Fakt}
{Teorema}
{}
{

\faktsituation {Misalkan $X$ suatu
\definitionsverweis {manifold Riemann}{}{.}
Maka
\definitionsverweis {koneksi Levi--Civita}{}{}
$\nabla$ memenuhi sifat-sifat berikut.}
\faktfolgerung {\aufzaehlungdrei{ $\nabla$
\definitionsverweis {linear}{}{.}
}{ $\nabla$ metrik, yakni

\mavergleichskettedisp
{\vergleichskette
{ D_V \left\langle  W  ,  Z  \right\rangle
}
{ =} { \left\langle \nabla_V W , Z  \right\rangle  + \left\langle  W  ,  \nabla_VZ  \right\rangle
}
{ } {
}
{ } {
}
{ } {
}
}
{}{}{.}
}{ $\nabla$ bebas torsi, yakni

\mavergleichskettedisp
{\vergleichskette
{ \nabla_V W- \nabla_W V
}
{ =} { [V,W]
}
{ } {
}
{ } {
}
{ } {
}
}
{}{}{.}
}}
\faktzusatz {}
\faktzusatz {}

}
{

Hal ini mengikuti dari
Lemma 26.11.

}



\inputfaktbeweis
{R^3/Riemannsche orientierte Fläche/Levi-Civita-Zusammenhang und induzierter Zusammenhang/Fakt}
{Lemma}
{}
{

\faktsituation {Misalkan

\mavergleichskette
{\vergleichskette
{ Y
}
{ \subseteq }{ \R^3
}
{  }{
}
{    }{
}
{   }{
}
}
{}{}{}
suatu permukaan yang dua kali terdiferensialkan dan dilengkapi oleh ruang sekitar $\R^3$ dengan
\definitionsverweis {struktur Riemann}{}{}
terinduksi.}
\faktfolgerung {Maka
\definitionsverweis {koneksi Levi--Civita}{}{}
pada $Y$ berimpit dengan apa yang
\zusatzklammer {melalui
\definitionsverweis {proyeksi ortogonal}{}{}
diberikan pada ruang-ruang tangen} {} {}
didefinisikan sebagai
\definitionsverweis {koneksi}{}{}
dari
Catatan 24.8.}
\faktzusatz {}
\faktzusatz {}

}
{

Kedua koneksi itu linear. Karena itu, cukup ditunjukkan secara lokal bahwa
\definitionsverweis {turunan vertikal}{}{}
terkait sama pada medan-medan basis. Misalkan

\mavergleichskette
{\vergleichskette
{ V
}
{ \subseteq }{ \R^2
}
{  }{
}
{    }{
}
{   }{
}
}
{}{}{}
\definitionsverweis {terbuka}{}{}
dan
\maabbdisp {\varphi} { V } { U
} {}
suatu parametrisasi lokal yang dua kali terdiferensialkan secara kontinu dari sebuah himpunan bagian terbuka

\mavergleichskette
{\vergleichskette
{ U
}
{ \subseteq }{ Y
}
{  }{
}
{    }{
}
{   }{
}
}
{}{}{.}
Misalkan $\partial_1,\partial_2$ adalah kedua medan arah konstan pada $V$. Untuk koneksi Levi--Civita berlaku

\mavergleichskettedisp
{\vergleichskette
{ \nabla_{\partial_i } \partial_j
}
{ =} { \sum_{k = 1}^2  \Gamma^k_{ij} \partial_k
}
{ } {
}
{ } {
}
{ } {
}
}
{}{}{,}
dengan simbol-simbol Christoffel yang didefinisikan berdasarkan

\mavergleichskettedisp
{\vergleichskette
{ g_{ij}
}
{ =} { \left\langle  \partial_i \varphi ,  \partial_j \varphi  \right\rangle
}
{ } {
}
{ } {
}
{ } {
}
}
{}{}{.}
Menurut
Lemma 26.5,
simbol-simbol Christoffel ini juga memenuhi

\mavergleichskettedisp
{\vergleichskette
{ \partial_i \partial_j \varphi
}
{ =} { \Gamma^1_{ij} \partial_1  \varphi +  \Gamma^2_{ij}  \partial_2  \varphi+   h_{ij} N
}
{ } {
}
{ } {
}
{ } {
}
}
{}{}{,}
dengan $N$ menyatakan
\definitionsverweis {medan normal satuan}{}{}
dan $h_{ij}$ merupakan entri-entri
\definitionsverweis {matriks fundamental kedua}{}{}
\zusatzklammer {semuanya dipandang pada $V$} {} {.}

Medan-medan vektor $\partial_i$ pada $V$ bersesuaian, pada

\mavergleichskette
{\vergleichskette
{ U
}
{ \subseteq }{ Y
}
{ \subseteq }{ \R^3
}
{    }{
}
{   }{
}
}
{}{}{}
dengan medan-medan vektor $\left(  \partial_i \varphi  \right)  \circ \varphi^{-1}$. Menurut definisi
\definitionsverweis {turunan vertikal}{}{}
untuk koneksi terbatas tersebut, kita harus meninjau

\mathdisp {U \stackrel{ \left(  \partial_i\varphi  \right) \circ \varphi^{-1} }{\longrightarrow} TU \stackrel{ T \left(  \left(  \partial_j \varphi  \right) \circ \varphi^{-1}    \right) }{\longrightarrow} TTU \stackrel{ \pi_{\text{vert} } } \longrightarrow TU} {  }

Ini menghasilkan proyeksi ortogonal
\mathl{\left(  \partial_i \partial_j \varphi  \right) \circ \varphi^{-1}}{}
ke bundel tangen pada $Y$, yang berimpit dengan
\mathl{\Gamma^1_{ij} \partial_1  \varphi +  \Gamma^2_{ij}  \partial_2  \varphi}{,}
dipandang di atas $U$
\zusatzklammer {lihat juga
Latihan 26.11} {} {.}

}

Pernyataan ini juga berlaku dengan cara yang sama untuk setiap submanifold tertutup

\mavergleichskette
{\vergleichskette
{ Y
}
{ \subseteq }{ \R^n
}
{  }{
}
{    }{
}
{   }{
}
}
{}{}{}
yang dilengkapi dengan struktur Riemann terinduksi.
